\documentclass[12pt]{article}
\usepackage{graphicx} % Required for inserting images
\usepackage{amsmath}
\usepackage{amssymb}
\usepackage{amsfonts}
\usepackage
[papersize={8.5in,11in},text={7.2in,9.5in}]
{geometry}

\usepackage{amsmath}
\newcommand{\bm}[1]{\boldsymbol{#1}}
\newcommand{\norm}[1]{\left \lVert #1 \right \rVert}

% Either use DeclareMathOperator (preferred):
\DeclareMathOperator*{\argmin}{arg\,min}
\DeclareMathOperator*{\argmax}{arg\,max}
\title{Investigating reduced-FOV perfusion FMRI using 
velocity selective inversion ASL with outer-volume suppression
}
\author{Yongli He, Jon-Fredrik Nielsen, Luis Hernandez-Garcia}
\date{}

\begin{document}

\maketitle
\section{Impact(40 words)}


\section{Synopsis(100 words)}
\subsection{Motivation}

\subsection{Goal}

\subsection{Approach}

\subsection{Results}


\newpage
\section{Introduction}
Velocity selective arterial spin labeling (VS-ASL) has been an attractive option 
for perfusion imaging because of its insensitivity to arterial transit time[cite],
which offers potential SNR gain and quantification bias reduction
compared to conventional spatial selective counterparts (e.g., PCASL[cite]).

Compared to BOLD fMRI, perfusion signal can quantify cerebral blood flow (CBF), 
providing higher spatial specificity to individual cortical column 
where GRE-BOLD tends to be confounded by venous contamination[Duong et al 2001].
This makes it attractive to laminar fMRI,
which to date is still limited in 7T and requires 
vascular contrasts such as VASO[Huber et al 2017].

Here we present reduced-FOV VSASL fMRI with 
tailored 2D outer-volume-suppresion (OVS) pulse. 
We hypothesize that OVS permits an accelerated readout with reduced under-sampling aliasing,
and reduces variance in the region of interest (ROI).
This work paves way for a long term goal of 
achieving laminar perfusion fMRI in 3T. 
\section{Methods}
\subsection{Sequence}
Each TR (Fig.~1a) comprises a $90^\circ$ presaturation, recovery delay D0, 
the velocity-selective label, two background-suppression (BGS) inversions,
a vascular crusher, an OVS prep train, and a stack-of-spirals (SoS) readout.

Labeling uses the FT-VSI pulse of Qin et al~\cite{qin2016velocity}
(Fig.~1b), with its velocity-encoding gradient scaled
to a cutoff velocity $V_c=2$~cm/s (15~mT/m). 
A BIR-8 velocity-selective saturation crusher (Fig.~1c) 
suppresses residual intravascular label with the same $V_c$. 
Two adiabatic sech inversions are inserted during post labeling delay (PLD)
as background suppresion (BGS).
We optimize the BGS timing D1-D3 as described in the next section.

The tailored 2D OVS pulse~\cite{he2026improved} (Fig.~1d) precedes every readout
excitation; because its suppression requires several repetitions to reach steady
state, five OVS-prep modules precede the readout train. The readout (Fig.~1e) is an
RF-spoiled SPGR SoS acquisition with a constant $10^\circ$ flip angle.

\subsection{Optimization of timing parameters}
D1 and D2 were chosen to null static CSF at readout~\cite{hu2026csf}, 
using a Bloch-free longitudinal-recovery model (Figure 2)
with $T_1$ = 4.3/1.2/0.8~s for CSF/GM/WM at 3~T[cite]. 
Each adiabatic inversion was assumed 90\% efficient, and the
FT-VSI label was assumed to invert static spins to $M_z=-0.81$; 
the velocity-selective modules were assumed not to otherwise attenuate static magnetization. 
D0 = 2~s and PLD = 1.4~s were fixed. 
We then minimized the residual $M_z$ of CSF, 
with GM and WM suppression as a secondary term,
over D1, D2 $\geq$ 20~ms by multi-start Nelder--Mead, 
subject to D3 $\geq$ 10~ms so
the crusher and OVS-prep train still fit inside the PLD, 
which caps D1\,+\,D2 at 1.307~s.
The optimum was D1 = 1.157~s, D2 = 0.149~s (D3 = 10~ms), 
predicting $M_z/M_0$ of
$-0.004$, $+0.22$ and $+0.33$ for CSF, GM and WM at readout (Fig.~2).
\subsection{FMRI experiments}


\section{Results}


\section{Discussion/Conclusion}

\section{Acknowledgments}




\bibliographystyle{unsrt}
\bibliography{myref}
\end{document}
